\begin{summarybox}
	\textbf{\textcolor{CSummary}{一句话核心}}

	正四面体的外接球半径 $R$ 与棱长 $a$ 和高 $h$ 满足
	\[
		R = \frac{\sqrt{6}}{4}a = \frac{3}{4}h,
	\]
	比例固定。

	\medskip
	\textbf{\textcolor{CSummary}{使用条件}}
	\begin{itemize}[leftmargin=*, itemsep=0.5em, topsep=0.4em]
		\item \textbf{条件 1（几何体限定）：} 几何体必须是所有棱长都相等的正四面体。
		\item \textbf{条件 2（已知元素）：} 已知棱长 $a$、高 $h$ 或外接球半径 $R$ 中任意一个，可求其他相关量。
		\item \textbf{条件 3（防止误用）：} 一般三棱锥或正三棱锥不能直接套用本公式；只有六条棱全部相等时才可使用。
	\end{itemize}

	\medskip
	\textbf{\textcolor{CSummary}{关键提醒}}
	\begin{itemize}[leftmargin=*, itemsep=0.5em, topsep=0.4em]
		\item \textbf{易错点（公式混淆）：} 不要将外接球半径公式
		      \[
			      R=\frac{\sqrt{6}}{4}a
		      \]
		      与高公式
		      \[
			      h=\frac{\sqrt{6}}{3}a
		      \]
		      混淆；两者系数不同。
		\item \textbf{易错点（球心位置）：} 外接球球心在高线上，距离底面为 $\frac{1}{4}h$，距离顶点为 $\frac{3}{4}h$，因此球心更靠近底面。
		\item \textbf{易错点（对称概念）：} 正四面体具有高度对称性，但不是中心对称图形；在正四面体中，外接球球心、内切球球心、重心和垂心重合。
		\item \textbf{检查项（条件确认）：} 使用公式前必须确认几何体为正四面体，否则公式不成立。
	\end{itemize}
\end{summarybox}